\begin{lemma}[Every member of the family is the jet of a $p$-harmonic function]
Let $1<p<2$, let $(\kappa,C_0)$ be a first-order estimate datum for $p$, and let
$z\in\mathcal{K}_p$.  Then
\begin{enumerate}
\item[(a)] $\bigl(z(x)\bigr)_1=0$ for every $x\in X_{1/2}$ with $x=0$ as a point of the plane;
\item[(b)] $\bigl\lvert\bigl(z(x)\bigr)_1\bigr\rvert\le1$ for every $x\in X_{1/2}$;
\item[(c)] there are $h\colon\mathrm{Plane}\to\mathbb{R}$ and
$G\colon\mathrm{Plane}\to\mathrm{Plane}$ such that $h$ is $p$-harmonic in $B_{1/2}$ with weak
gradient $G$, such that $h(x)=\bigl(z(x)\bigr)_1$ and $G(x)=\bigl(z(x)\bigr)_2$ for every
$x\in X_{1/2}$, and such that
\[
  G(y)=Dh(y)\qquad\text{for every }y\in B_{1/2},
\]
where $Dh$ is the classical gradient of $h$.
\end{enumerate}
\end{lemma}
\begin{proof}
By \noderef{CompactJetFamily}, $\mathcal{K}_p=\overline{\mathcal{G}_p}$, so
$z\in\overline{\mathcal{G}_p}$.

\medskip\noindent\emph{Step 0: inclusions of balls.}
For $x\in X_{1/2}$ we have $\lvert x\rvert\le\tfrac12<\tfrac34<1$, so
\begin{equation}\label{eq:cjp-incl}
  X_{1/2}\subseteq B_{3/4}\subseteq B_1 ,
\end{equation}
and for $y\in B_{1/2}$ we have $\lvert y\rvert<\tfrac12<\tfrac34$, so
\begin{equation}\label{eq:cjp-incl2}
  B_{1/2}\subseteq B_{3/4},\qquad B_{1/2}\subseteq\overline{B_{1/2}}=X_{1/2}\subseteq B_1 .
\end{equation}

\medskip\noindent\emph{Step 1: a sequence of generators converging to $z$.}
The space $X_{1/2}\to^{b}(\mathbb{R}\times\mathrm{Plane})$ is a metric space under the supremum
distance, by clause (1) of the background proposition on bounded continuous functions.  Clause
(3) of that proposition --- in a metric space every point of the closure of a set is the limit
of a sequence of its elements --- therefore gives a sequence $(z_n)_{n\in\mathbb{N}}$ in
$\mathcal{G}_p$ with $z_n\to z$ in the supremum metric.  By clause (1) again, $z_n\to z$
uniformly on $X_{1/2}$, and in particular $z_n(x)\to z(x)$ for every $x\in X_{1/2}$; by clause
(2), that pointwise convergence holds in each component separately.

\medskip\noindent\emph{Step 2: unpacking the generators.}
Fix $n$.  Since $z_n\in\mathcal{G}_p$, \noderef{CompactJetFamily} provides
$v_n\colon\mathrm{Plane}\to\mathbb{R}$ and $g_n\colon\mathrm{Plane}\to\mathrm{Plane}$ with
\begin{enumerate}
\item[(i)] $v_n$ is $p$-harmonic in $B_1$ with weak gradient $g_n$;
\item[(ii)] $\lvert v_n(y)\rvert\le1$ for every $y\in B_1$;
\item[(iii)] $v_n(0)=0$;
\item[(iv)] $z_n(x)=\bigl(v_n(x),Dv_n(x)\bigr)$ for every $x\in X_{1/2}$.
\end{enumerate}

\medskip\noindent\emph{Step 3: parts (a) and (b).}
Let $x\in X_{1/2}$ with $x=0$ as a point of the plane.  By (iv) and (iii),
$\bigl(z_n(x)\bigr)_1=v_n(0)=0$ for every $n$.  By Step 1 the first components converge,
$\bigl(z_n(x)\bigr)_1\to\bigl(z(x)\bigr)_1$, so clause (6) of the
bounded-continuous-functions proposition --- a pointwise limit of a constant sequence is that
constant --- gives $\bigl(z(x)\bigr)_1=0$.  This is part (a).

Let now $x\in X_{1/2}$ be arbitrary.  By \eqref{eq:cjp-incl} we have $x\in B_1$, so (ii) and
(iv) give $\bigl\lvert\bigl(z_n(x)\bigr)_1\bigr\rvert=\lvert v_n(x)\rvert\le1$, that is
$-1\le\bigl(z_n(x)\bigr)_1\le1$, for every $n$.  Applying clause (6) of the same proposition to
the convergent sequence $\bigl(z_n(x)\bigr)_1\to\bigl(z(x)\bigr)_1$ --- once as stated, and once
to the negated sequence for the lower bound --- gives $-1\le\bigl(z(x)\bigr)_1\le1$, that is
part (b).

\medskip\noindent\emph{Step 4: the extension off the closed half-ball.}
Define $h\colon\mathrm{Plane}\to\mathbb{R}$ and $G\colon\mathrm{Plane}\to\mathrm{Plane}$ by
\[
  h(y):=\begin{cases}\bigl(z(y)\bigr)_1,&y\in\overline{B_{1/2}},\\ 0,&\text{otherwise,}\end{cases}
  \qquad
  G(y):=\begin{cases}\bigl(z(y)\bigr)_2,&y\in\overline{B_{1/2}},\\ 0,&\text{otherwise,}\end{cases}
\]
which is legitimate because $\overline{B_{1/2}}$ is exactly the set of points of $X_{1/2}$, so
$z(y)$ is defined for $y\in\overline{B_{1/2}}$.  By construction
\begin{equation}\label{eq:cjp-agree}
  h(x)=\bigl(z(x)\bigr)_1\quad\text{and}\quad G(x)=\bigl(z(x)\bigr)_2
  \qquad\text{for every }x\in X_{1/2},
\end{equation}
which is the second half of part (c).  It remains to prove that $h$ is $p$-harmonic in $B_{1/2}$
with weak gradient $G$.

\medskip\noindent\emph{Step 5: uniform convergence of the extensions.}
Let $\varepsilon>0$.  By Step 1 there is $N$ such that
$d\bigl(z_n(x),z(x)\bigr)\le\varepsilon$ for every $n\ge N$ and every $x\in X_{1/2}$.  Fix
$n\ge N$ and $y\in\overline{B_{1/2}}$, so that $y$ is a point of $X_{1/2}$.  By (iv),
\eqref{eq:cjp-agree} and clause (2) of the bounded-continuous-functions proposition --- the
distance between either component is at most the distance between the pairs ---
\[
  \lvert v_n(y)-h(y)\rvert
  =\bigl\lvert\bigl(z_n(y)\bigr)_1-\bigl(z(y)\bigr)_1\bigr\rvert
  \le d\bigl(z_n(y),z(y)\bigr)\le\varepsilon ,
\]
\[
  \lvert Dv_n(y)-G(y)\rvert
  =\bigl\lvert\bigl(z_n(y)\bigr)_2-\bigl(z(y)\bigr)_2\bigr\rvert
  \le d\bigl(z_n(y),z(y)\bigr)\le\varepsilon .
\]
Since $\varepsilon>0$ was arbitrary and $N$ does not depend on $y$, this says
\begin{equation}\label{eq:cjp-unif}
  v_n\to h\quad\text{and}\quad Dv_n\to G\qquad\text{uniformly on }\overline{B_{1/2}} .
\end{equation}
In particular, every compact $K\subseteq B_{1/2}$ satisfies
$K\subseteq B_{1/2}\subseteq\overline{B_{1/2}}$, so $v_n\to h$ and $Dv_n\to G$ uniformly on
every compact subset of $B_{1/2}$.

\medskip\noindent\emph{Step 6: each $v_n$ is $p$-harmonic in $B_{1/2}$ with its classical
gradient.}
Fix $n$.

\emph{(a) Restriction.}  The ball $B_{1/2}$ is open and $B_{1/2}\subseteq B_1$ by
\eqref{eq:cjp-incl2}, so part (a) of \noderef{PHarmRestrict}, applied to the datum (i), gives
that $v_n$ is $p$-harmonic in $B_{1/2}$ with weak gradient $g_n$.

\emph{(b) $C^1$ regularity.}  Apply \noderef{PHarmJetEquicontinuity} with $x_0:=0$,
$R:=\tfrac12$ and $M:=1$.  Its hypotheses hold: $1<p<2$; $R>0$ and $M>0$; $v_n$ is $p$-harmonic
in $B_{2R}(0)=B_1$ with weak gradient $g_n$, which is (i), since $2\cdot\tfrac12=1$; and
$\lvert v_n(y)\rvert\le1=M$ for $y\in B_1$, which is (ii).  Since $3R/2=\tfrac34$, its first
conclusion is that $v_n$ is of class $C^1$ on $B_{3/4}$.  As $B_{1/2}\subseteq B_{3/4}$ by
\eqref{eq:cjp-incl2}, clause (4) of the background proposition on locality of $C^1$ regularity
gives that $v_n$ is of class $C^1$ on $B_{1/2}$.

\emph{(c) The bridge.}  Now \noderef{PHarmWeakGradClassical} applies with $\Omega:=B_{1/2}$,
$u:=v_n$ and $g:=g_n$: its hypotheses are $p>1$, which follows from $1<p$; that $v_n$ is
$p$-harmonic in $B_{1/2}$ with weak gradient $g_n$, which is step (a); and that $v_n$ is of
class $C^1$ on $B_{1/2}$, which is step (b).  Its part (a) gives
\begin{equation}\label{eq:cjp-classical}
  v_n\ \text{is}\ p\text{-harmonic in}\ B_{1/2}\ \text{with weak gradient}\ Dv_n .
\end{equation}
This is the step that the generating condition of \noderef{CompactJetFamily} leaves open: that
condition supplies only the arbitrary weak gradient $g_n$, whereas what converges in Step 5 is
the classical gradient $Dv_n$, and the two agree only almost everywhere.

\medskip\noindent\emph{Step 7: limit passage.}
Apply \noderef{PHarmLimit} with $\Omega:=B_{1/2}$, with the sequences $n\mapsto v_n$ and
$n\mapsto Dv_n$, and with the limits $h$ and $G$.  Its hypotheses hold: $1<p<2$; for every $n$
the function $v_n$ is $p$-harmonic in $B_{1/2}$ with weak gradient $Dv_n$, which is
\eqref{eq:cjp-classical}; and for every compact $K\subseteq B_{1/2}$ one has $v_n\to h$ and
$Dv_n\to G$ uniformly on $K$, which is the last sentence of Step 5.  Its conclusion is that $h$
is $p$-harmonic in $B_{1/2}$ with weak gradient $G$.

\medskip\noindent\emph{Step 8: the weak gradient is the classical gradient of $h$ on
$B_{1/2}$.}
Apply \noderef{PHarmClassicalGradient} with the open set $\Omega:=B_{1/2}$, the sequences
$n\mapsto v_n$ and $n\mapsto Dv_n$, and the limits $h$ and $G$.  Its three hypotheses hold.
First, for every $n$ and every $y\in B_{1/2}$ the function $v_n$ is differentiable at $y$ with
classical gradient $Dv_n(y)$: indeed $v_n$ is of class $C^{1}$ on $B_{1/2}$ by Step 6(b), so
clause (2) of the background proposition on locality of $C^{1}$ regularity makes it
differentiable at every point of that open set, and its classical gradient there is by
definition $Dv_n$.  Second, for every compact $K\subseteq B_{1/2}$ we have $Dv_n\to G$
uniformly on $K$, which is the last sentence of Step 5.  Third, for every $y\in B_{1/2}$ we have
$v_n(y)\to h(y)$, since by the same sentence the convergence $v_n\to h$ is uniform on the
compact set $\{y\}\subseteq B_{1/2}$, hence pointwise there.  The conclusion of
\noderef{PHarmClassicalGradient} is that for every $y\in B_{1/2}$ the function $h$ is
differentiable at $y$ with classical gradient $G(y)$, and that $Dh(y)=G(y)$.  This is the last
assertion of part (c).

Note that this step genuinely needs the approximating sequence, which exists only inside this
proof: a consumer holding only $z$, $h$ and $G$ knows that $G$ is \emph{a} weak gradient of $h$,
and weak gradients are determined only up to a null set, so the pointwise identity could not be
recovered downstream.  That is why it is recorded in the statement.

\medskip\noindent Together with \eqref{eq:cjp-agree} and Step 8 this establishes part (c), and
parts (a) and (b) were proved in Step 3.
\end{proof}
